\begin{center}
    \textbf{ВВЕДЕНИЕ}
\end{center}
\addcontentsline{toc}{chapter}{ВВЕДЕНИЕ}

Обеспечение безопасного доступа к файлам в операционной системе Linux является актуальной задачей. При работе с файлами в Linux необходимо обеспечивать конфиденциальность данных и предоставлять защиту от вредоносных действий.

Целью данной курсовой работы является разработка загружаемого модуля ядра для ОС Linux, позволяющего скрывать файлы или запрещать их изменение, чтение и удаление. Необходимо предусмотреть возможность ввода пароля для отображения файлов или разрешения операций над ними. Предоставить пользователю возможность задавать список таких файлов.

Для решения поставленной задачи необходимо:
\begin{itemize}
	\item[---] проанализировать возможности перехвата функций ядра Linux;
	\item[---] выбрать системные вызовы, которые необходимо перехватить;
	\item[---] разработать алгоритм перехвата, алгоритмы hook–функций и структуру программного обеспечения;
	\item[---] реализовать программное обеспечение;
	\item[---] исследовать работу ПО.
\end{itemize}

 